Parameter estimation for nonlinear dynamical systems from time-series observations is an inverse problem with broad scientific applications~\cite{tarantola2005inverse}. 
While parameter estimation is inherently ill-posed, practical constraints on data collection severely exacerbate this difficulty.
In practical settings, observations are frequently restricted:
the system state is only \emph{partially} observable, capturing a limited subset of the variables, and the temporal sampling intervals can become highly \emph{sparse}.
These limitations deprive the inverse mapping of critical information, frequently driving the system into non-identifiable regimes where distinct parameter values induce essentially indistinguishable observed trajectories.

A common approach is to treat inference as either (i) per-instance numerical optimization over parameters (and possibly latent states) or (ii) direct regression from observations to parameters using a machine learning model. 
While optimization-based methods are mathematically principled, they can be computationally expensive and sensitive to initialization in sparse regimes, and they may converge to different solutions when the inverse problem admits multiple nearly equivalent explanations. 
Regression-based approaches amortize inference cost by learning a direct map from observations to parameters. They are simple and, as our experiments confirm, a strong baseline; what they do not do is exploit the dynamical structure of the forward model or form any explicit estimate of the unobserved state. This gap motivates methods that reconstruct the latent trajectory as part of inference, on the intuition that using the hidden state should help --- an intuition whose limits this paper makes precise.

In this work, we adopt a fixed-point inference formulation that makes the inference dynamics explicit.
We represent sparse observations $\mathbf{x}_{\mathrm{obs}}$  and define a learned update operator $ \mathcal{T}(\theta;\mathbf{x}_{\mathrm{obs}})=P_{\psi}\left(\mathbf{x}_{\mathrm{obs}},H_{\phi}(\mathbf{x}_{\mathrm{obs}},\theta)\right)$, where $H_{\phi}$ is a hidden-state estimator and $P_{\psi}$ is a parameter estimator. 
Inference is performed by the fixed-point iteration $\theta^{k+1}=\mathcal{T}(\theta^{k};\mathbf{x}_{\mathrm{obs}})$, and the final prediction is the fixed point $\theta^*(\mathbf{x}_{\mathrm{obs}})$. 
The key design objective is contractivity of $\mathcal{T}(\cdot;\mathbf{x}_{\mathrm{obs}})$ for each fixed $\mathbf{x}_{\mathrm{obs}}$, which guarantees global convergence to a unique fixed point with a geometric rate.

Crucially, convergence does not imply correctness, and analyzing \emph{what} the fixed point is turns out to yield a characterization of a fundamental limitation rather than a guarantee of accuracy. Because the fixed point is a deterministic function of $\mathbf{x}_{\mathrm{obs}}$, it cannot be more accurate than the conditional mean $\mathbb{E}[\theta\mid\mathbf{x}_{\mathrm{obs}}]$; in particular, reconstructing the hidden state --- the feature that motivates the decoupling --- confers no information advantage at inference, since the reconstructed state is itself an image of $\mathbf{x}_{\mathrm{obs}}$. The deviation of the fixed point from this conditional-mean ceiling is controlled by a one-step residual $\varepsilon(\mathbf{x}_{\mathrm{obs}})=\lVert \mathcal{T}(\theta^{\mathrm{true}};\mathbf{x}_{\mathrm{obs}})-\theta^{\mathrm{true}}\rVert$, which we bound and decompose into a structural part and an operator-approximation part.

The consequence is sharpest under non-identifiability, where the conditional mean is not merely inaccurate: for a curved solution fiber it lies strictly \emph{off} the manifold of parameters consistent with the observation, so the squared-loss estimate is not even physically admissible. This regression-to-the-mean and joint-distribution collapse persists no matter how well the iteration converges; it reflects the information content of the observations and the choice of supervised loss, not a numerical artifact. Taken together, our results say that a decoupled deterministic point estimator, however well-motivated and however stably it converges, is bounded by the (off-fiber) conditional mean --- a limitation that motivates lifting inference from point estimates to distributions.

Our main contributions are:
\begin{itemize}
    \item We formulate parameter estimation from sparse partial observations as a fixed-point problem in parameter space, defined by an update operator $\theta^{(i+1)}=\mathcal{T}(\theta^{(i)};\mathbf{x}_{\mathrm{obs}})$ composed of a hidden-state estimator and a parameter estimator.
    \item We provide a sufficient condition for $\mathcal{T}(\cdot;\mathbf{x}_{\mathrm{obs}})$ to be contractive and hence globally convergent to a unique fixed point with a geometric rate, and we show how to enforce the required Lipschitz bounds in practice (e.g., via spectral normalization).
    \item  We separate convergence from correctness and characterize the fixed point: it is a function of $\mathbf{x}_{\mathrm{obs}}$ and therefore cannot beat the conditional-mean ceiling $\mathbb{E}[\theta\mid\mathbf{x}_{\mathrm{obs}}]$ (reconstructing the hidden state adds no information at inference); its deviation from that ceiling is bounded by a one-step residual $\varepsilon(\mathbf{x}_{\mathrm{obs}})$, which we decompose into a structural and an operator-approximation part.
    \item We show that under non-identifiability the conditional mean is itself displaced off the physical solution fiber (a Jensen/extreme-point argument), so squared-loss point estimation yields a physically inadmissible estimate; this regression-to-the-mean and joint-distribution collapse persists regardless of convergence.
    \item On epidemiological (SIR), ecological (Lotka--Volterra), and physiological (OGTT) systems, whose identifiability we predict a priori by a Hermann--Krener observability analysis, we confirm each prediction empirically --- including that a direct regressor reaches the conditional-mean ceiling more closely than the decoupled iteration, isolating the cost of the reconstruction detour.
\end{itemize}
